\section{Considerações finais}\label{cap:conclusao}
% orçamento: 350 palavras

Partimos de uma pergunta de um trabalho de conclusão de curso: uma soma de círculos girando consegue descrever a órbita de um planeta, e os círculos de Copérnico têm a ver com essa soma? Para responder construímos um simulador que integra a gravitação mútua do Sol e dos oito planetas a partir de dados do JPL, com um integrador simplético de quarta ordem. Nas cinco datas testadas, entre $-100$ e $+100$ anos, ele reproduz a efeméride de referência com erro abaixo de 0,5\,\% da distância ao Sol, sendo Mercúrio o pior caso, e a energia oscila sem tendência por mil anos.

Sobre essas órbitas, a série de Fourier complexa parametrizada no tempo tem, até primeira ordem na excentricidade, a estrutura da construção do \emph{Commentariolus}: um deslocamento fixo de $\frac32 ae$, um deferente de raio $a$ e um epiciclo de raio $\frac12 ae$ girando duas vezes por período. Os coeficientes têm forma fechada em funções de Bessel, conferida numericamente até o arredondamento da máquina. Os números de Copérnico para Marte, Júpiter e Saturno concordam com esses valores em poucos por cento, e a razão de três para um entre os seus dois epiciclos é a que a série impõe.

Para saber se essa concordância é própria de Copérnico, usamos como controle o equante de Ptolomeu. Com parâmetros modernos, o equante e a série erram a posição de Marte em 0,007\,UA e o modelo de Copérnico em 0,013\,UA, os três em ordem $e^2$. O mecanismo de Copérnico tem o mesmo conteúdo de primeira ordem que o equante, e a série de Fourier explica por que os dois funcionam. Com os números históricos, nenhum dos dois erra sistematicamente menos em longitude. Pôr a origem no Sol médio, com as linhas das apsides modernas, piora Marte, Júpiter e Saturno e não explica a excentricidade de Vênus.

Em termos de física, o trabalho liga as duas descrições do céu: Kepler e Newton explicam por que os epiciclos funcionavam, e a série de Fourier mostra como uma descrição se escreve na outra.
